\section{Proofs and numerical implementation}\label{app:proofs}\label{app:formal_model}

This appendix proves the propositions of Section~\ref{sec:model}, documents the numerical solver, and reports the reduced-form benchmark referenced in Section~\ref{sec:identification}. The specification is the one-class representative-holder model of Section~\ref{sec:model}; the multi-class generalisation reuses the same fixed-point machinery and is taken up in the internet appendix.

\subsection{Proofs}\label{app:proof_block}

\begin{proof}[Proof of Proposition~\ref{prop:existence}]
The best-response operator $T$ maps a candidate threshold $\tilde s$ to the threshold that makes the marginal holder indifferent given that all other holders use $\tilde s$. Under the stated regularity, $T$ is continuous and bounded; the contraction hypothesis, which we verify numerically (the maximum row sum of the discrete Jacobian of $T$ is approximately $0.008$ at the calibrated parameter vector), implies a unique fixed point by Banach's theorem. The fixed point of $T$ coincides with the zero of $F$ in \eqref{eq:F} by construction. \qed
\end{proof}

\begin{proof}[Proof of Proposition~\ref{prop:reserves}]
A higher $R$ reduces $U(D)$ by \eqref{eq:U}, raising $p^{sec}$ in \eqref{eq:psec} and hence $V_R$ at any $D$. The set of $\theta$ for which the marginal holder prefers to redeem expands, shifting the equilibrium threshold $s^*$ up, so the aggregate $D^*(\theta)$ rises pointwise in $\theta$ over the support where the strategy is binding; integration over the prior preserves the sign. The economic content is that better reserves anticipate a higher secondary price, which encourages holders to redeem rather than ride out a temporary discount. \qed
\end{proof}

\begin{proof}[Proof of Proposition~\ref{prop:convenience}]
A higher $\omega$ raises $V_H$ in \eqref{eq:VH} at any $\theta$, so the set of $\theta$ for which the marginal holder prefers to redeem contracts, the equilibrium threshold $s^*$ falls, and the aggregate $D^*(\theta)$ falls pointwise in $\theta$. The economic content is that a higher payment-use value of holding the stablecoin lowers the equilibrium conversion mass: a useful stablecoin retains holders even at marginally lower fundamentals. \qed
\end{proof}

\subsection{Numerical implementation}\label{app:numerical}

The indifference equation \eqref{eq:F} is solved with \texttt{scipy.optimize.root} (Krylov method); the expectation over $\theta$ uses Gauss-Hermite quadrature with 15 nodes against the posterior $\mathcal{N}(s^*, \sigma^2)$. Each Monte Carlo draw fixes a scenario, draws the latent fundamental from the prior, and solves the fixed point at the realised $\theta$; the vectorised solver runs at approximately 29 milliseconds per draw on a single core. The realised secondary price and redemption mass at each $\theta$ feed into the moments and the risk-component decomposition of Section~\ref{sub:what_it_delivers}.

\subsection{A reduced-form benchmark}\label{app:decomposition}

A reduced-form alternative would price only counterparty and run, as a run-probability model does: $\Pi_{\text{red}} = \ell_C + p_s\,\ell_{D,s}$, with no baseline-regime term, since in normal times there is no run \citep{klee_etal_2026_fragility}. Granting it our own stress cost $\ell_{D,s}$ (conservative, as a market-observed loss given run would be the intermediary's $0.902$ trough (internet appendix), not the holder's), $\Pi_{\text{red}}$ recovers $22.6$ of USDC's $48.9$ and $60.3$ of USDT's $71.2$ basis points. It is adequate for counterparty-dominated USDT ($85\%$ of the expected loss) but misses $54\%$ of USDC's: the baseline-regime depeg $\ell_{D,b}$ a run-probability framework cannot see. The asymmetry is the point: one recipe prices USDT roughly right and USDC roughly wrong, and it is what the per-issuer supervisory reading of Section~\ref{sub:policy_implications} rests on.
